\section{Experiment Construction}
\label{sec:experiment_construction}

\subsection{From Observational Data to Simulated Experiments}
\label{subsec:obs_to_exp}

Neither REES46 nor UCI Online Retail is an A/B test.
Both are observational logs of user activity on an e-commerce
platform.
We construct \emph{simulated experiment windows} by slicing
the continuous activity stream into fixed-length time intervals
and treating each interval as if it were a separate experiment.

\paragraph{Rationale.}
The prediction task studied in the paper is:
\emph{given $\Dpilot$ days of pilot data, forecast the number
of new users $\Utrue$ (and total triggers $\Ttrue$) that will
appear in the subsequent $\Dfollow$ days.}
This task depends only on the temporal structure of user arrivals
and engagement---not on whether the data was generated by a
randomized experiment.
An A/B test arm and an observational time window produce
identical data structures: a matrix
$A \in \mathbb{N}^{D \times N}$ of daily trigger counts per user.
The model sees the first $\Dpilot$ columns and predicts the rest.

\paragraph{Key assumption.}
The validity of this approach requires \emph{local stationarity}:
within each 28-day window, the user arrival process and
engagement dynamics are approximately stationary.
This is plausible for short windows on a large platform
(daily fluctuations average out), though it may be violated
during sharp seasonal transitions (e.g., Black Friday).
We do not assume stationarity \emph{across} windows---different
windows may have different traffic levels, user populations,
and engagement patterns.
This is a feature, not a limitation: it means our evaluation
spans a range of conditions, from low-traffic periods to
holiday peaks.

\paragraph{Connection to A/B testing.}
In a real A/B test, each treatment arm generates a stream of
user activity over the experiment's duration.
If the treatment does not affect the \emph{rate} at which users
arrive (only the outcome metric, e.g., conversion rate), then
the arrival process in each arm is identical to the observational
process.
Our simulated experiments therefore evaluate the model under
conditions that match the intended use case: forecasting
participation from early pilot data to support duration and
stopping decisions.

A key advantage of the REES46 dataset over the ASOS data is its
temporal extent: 7~months ($\sim$210~days) of continuous activity
allow us to construct many independent experiment windows from a
single data source.
This transforms the evaluation from a single-experiment anecdote
into a systematic study across diverse traffic conditions.

\subsection{Windowing Strategy}
\label{subsec:windowing}

\paragraph{Non-overlapping windows.}
We partition the full time span into non-overlapping windows of
$D = 28$~days.
With $\sim$210~days of data, this yields 7~experiments.
Within each window, the first $\Dpilot = 7$~days constitute the
pilot period and the remaining $\Dfollow = 21$~days the follow-up.

\paragraph{Sliding windows.}
To increase the number of experiments and smooth over
boundary effects, we also construct sliding windows with a
stride of 3~days.
This yields approximately 60~overlapping experiments,
providing a richer basis for calibration studies.
We use window sizes of $D \in \{14, 21, 28\}$ days with
corresponding pilot periods $\Dpilot \in \{5, 7\}$ to
assess robustness across different time horizons.

\paragraph{Rolling windows with varying follow-up horizons.}
{\color{purple}To assess robustness across different extrapolation horizons,
we construct rolling experiment windows with follow-up lengths
$k \in \{21, 50, 100\}$ days.
Each window has a fixed pilot of $\Dpilot = 7$ days and a follow-up
of $\Dfollow = k$ days, giving a total window length of $\Dpilot + k$.
Windows are placed at every possible start day (stride~1),
yielding approximately $D_{\mathrm{total}} - \Dpilot - k + 1$
experiments per value of $k$.
This produces 190 experiments for $k=21$, 161 for $k=50$, and
112 for $k=100$ on the 7-month REES46 dataset.}

\subsection{Ground Truth}
\label{subsec:ground_truth}

For each experiment window, we compute the ground-truth quantities
that our models aim to predict:
\begin{itemize}[leftmargin=*]
  \item $\Npilot$: the number of distinct users active during the
        pilot period (days $1, \ldots, \Dpilot$);
  \item $\Utrue = N_{\mathrm{total}} - \Npilot$: the number of
        \emph{new} users who first appear during the follow-up
        period (days $\Dpilot + 1, \ldots, D$);
  \item $\Ttrue$: the total number of triggers recorded during
        the follow-up period, summed across all users.
\end{itemize}

These are the estimands defined in Section~3 of the main paper.
The ability to compute exact ground truth for each window---rather
than relying on a single held-out period---is what enables the
calibration study requested by the reviewers.

\subsection{Experiment Summary}

\paragraph{REES46.}
The non-overlapping 28-day windowing produces 7~experiments
spanning a range of traffic levels, from quieter periods
(e.g., late March) to high-activity periods (e.g., November,
the holiday shopping season).
Pilot sizes $\Npilot$ range from several thousand to tens of
thousands of users, and the number of new users $\Utrue$ in
the follow-up is typically 2--5$\times$ the pilot size,
reflecting the rapid influx characteristic of eCommerce platforms.

\paragraph{UCI Online Retail.}
The 374-day span yields 13~non-overlapping 28-day windows.
Pilot sizes range from $\Npilot = 34$ (a low-activity period
in early January) to $\Npilot = 474$ (the holiday season),
providing a natural test of model robustness across varying
data sparsity.
Table~\ref{tab:uci_experiments_summary} in the supplementary
material provides the full experiment-by-experiment breakdown.

\subsection{Data Flow Diagram}

The complete pipeline from raw data to model evaluation is:
\begin{center}
\begin{tabular}{ccccc}
Raw CSVs & $\xrightarrow{\text{chunk + aggregate}}$ &
User--day table & $\xrightarrow{\text{window}}$ &
Experiment matrices \\
(47\,GB) & & ($\sim$50\,MB) & & ($\sim$1\,MB each)
\end{tabular}
\end{center}
\noindent
Each experiment matrix $A \in \mathbb{N}^{D \times N}$ is stored
as a NumPy array and passed directly to the model fitting code.
The raw CSVs are deleted after preprocessing; only the compact
metadata and matrices are retained.
